\chapter{Perimetre retenu : du signal legacy au dashboard}

\section{Objectif de cette version du rapport}
Cette version du memoire couvre uniquement la partie amont de la chaine
quantitative : production de signaux techniques et restitution de ces signaux
dans le tableau de bord.

Le perimetre est volontairement limite a :
\begin{enumerate}
\item la page legacy originale \texttt{/signals} (mode \texttt{view=legacy}),
\item la methodologie de detection et de filtrage des signaux,
\item la representation des resultats sur le dashboard V1.
\end{enumerate}

Le perimetre exclut la construction complete de strategie et le backtest final.

\section{Pourquoi ce cadrage}
Ce cadrage correspond a la phase la plus mature du projet:
\begin{itemize}
\item page signal legacy fonctionnelle avec navigation multi-niveaux ;
\item pipeline de signal OOS documente et expose via API ;
\item dashboard V1 consommant les scores agreges par horizon.
\end{itemize}

\section{Flux fonctionnel cible}
Le flux fonctionnel analyse dans ce rapport est :
\[
\text{OHLCV canonique} \rightarrow \text{Signal Engine legacy}
\rightarrow \text{Scores par famille} \rightarrow \text{Dashboard V1}
\]

Ce flux est la brique decisionnelle minimale: il transforme des donnees de
marche en un signal interpretable sans passer encore par la phase
strategie/backtest.

